\chapter{토크나이저 만들기}
\label{ch:tokenizer}

이 장에서는 텍스트를 숫자 시퀀스로 변환하는 \keyterm{토크나이저(tokenizer)}를 직접 구현한다. 두 가지 주요 알고리즘인 BPE(Byte Pair Encoding)와 Unigram을 다루며, 각각의 학습, 인코딩, 디코딩 과정을 코드로 살펴본다.

\section{왜 토크나이저가 필요한가}

신경망은 숫자만 이해한다. ``Hello, world!''라는 문자열을 그대로 모델에 넣을 수 없다. 텍스트를 정수 시퀀스로 변환하는 첫 단계가 필요한데, 이를 \keyterm{토크나이저}가 담당한다. 가장 단순한 방법은 ``글자 단위 토크나이저''로, 모든 문자를 별도의 토큰으로 취급하는 것이다. 하지만 이는 시퀀스가 너무 길어지고 단어의 의미를 잡기 어렵다.

반대 극단은 ``단어 단위 토크나이저''로, 공백으로 분리된 단어를 하나의 토큰으로 취급한다. 하지만 이는 어휘 크기가 폭발하고(영어만 해도 수십만 단어), 학습하지 않은 단어(out-of-vocabulary, OOV)를 처리하지 못한다.

현대 LLM은 이 두 극단의 중간인 \keyterm{서브워드(subword)} 토크나이저를 사용한다. 자주 등장하는 단어는 그대로 하나의 토큰으로, 희귀한 단어는 더 작은 서브워드로 분할한다. ``unhappiness''는 ``un'' + ``happiness'' 또는 ``un'' + ``happi'' + ``ness''로 쪼개질 수 있다. 이 방식은 어휘 크기를 적절히 유지하면서 OOV 문제를 해결한다.

\subsection{세 가지 접근 비교}

\begin{center}
\begin{tabularx}{\textwidth}{lXXX}
\toprule
\textbf{방식} & \textbf{어휘 크기} & \textbf{OOV 처리} & \textbf{시퀀스 길이} \\
\midrule
글자 단위 & 작음 ($\sim$100) & 없음 & 매우 김 \\
단어 단위 & 매우 큼 ($\sim$100K) & 처리 못함 & 짧음 \\
서브워드 & 중간 ($\sim$10K$\sim$50K) & 부분단어로 분해 & 중간 \\
\bottomrule
\end{tabularx}
\end{center}

\section{특수 토큰}

모든 토크나이저는 몇 가지 \keyterm{특수 토큰(special token)}을 가진다. 우리의 구현에서는 네 가지를 사용한다.

\begin{center}
\begin{tabularx}{\textwidth}{llX}
\toprule
\textbf{토큰} & \textbf{ID} & \textbf{용도} \\
\midrule
\code{<pad>} & 0 & 배치 내 시퀀스 길이를 맞추기 위한 패딩 \\
\code{<bos>} & 1 & 시퀀스의 시작을 표시 (Beginning Of Sequence) \\
\code{<eos>} & 2 & 시퀀스의 끝을 표시 (End Of Sequence) \\
\code{<unk>} & 3 & 어휘에 없는 토큰 (Unknown) \\
\bottomrule
\end{tabularx}
\end{center}

이 토큰들은 인덱스가 고정되어 있어 모델이 특별히 취급할 수 있다. 예를 들어 생성 시 \code{<eos>}가 나오면 추론을 멈추고, \code{<pad>} 위치는 loss 계산에서 무시한다.

\begin{warnbox}[특수 토큰 ID 고정의 이유]
특수 토큰의 ID를 0$\sim$3으로 고정하면 모델 코드가 단순해진다. ``\code{if token\_id == 0: mask\_out}'' 같은 하드코딩이 가능하다. 또한 체크포인트 호환성이 보장된다: 토크나이저를 다시 학습해도 특수 토큰 ID는 항상 같으므로 모델 가중치를 그대로 재사용할 수 있다.
\end{warnbox}

\section{베이스 클래스}

모든 토크나이저가 따라야 할 인터페이스를 추상 클래스로 정의한다.

\begin{lstlisting}[language=Python]
from abc import ABC, abstractmethod
from dataclasses import dataclass

@dataclass
class SpecialTokens:
    pad: str = "<pad>"   # ID 0
    bos: str = "<bos>"   # ID 1
    eos: str = "<eos>"   # ID 2
    unk: str = "<unk>"   # ID 3

    @property
    def tokens(self) -> list[str]:
        return [self.pad, self.bos, self.eos, self.unk]

    @property
    def pad_id(self) -> int: return 0
    @property
    def bos_id(self) -> int: return 1
    @property
    def eos_id(self) -> int: return 2
    @property
    def unk_id(self) -> int: return 3


class Tokenizer(ABC):
    def __init__(self, special_tokens: SpecialTokens | None = None):
        self.special = special_tokens or SpecialTokens()
        self.id_to_token: list[str] = list(self.special.tokens)
        self.token_to_id: dict[str, int] = {t: i for i, t in enumerate(self.id_to_token)}

    @abstractmethod
    def _train(self, corpus: list[str], vocab_size: int) -> None:
        """어휘 학습 (자식 클래스에서 구현)"""

    @abstractmethod
    def _tokenize(self, text: str) -> list[str]:
        """텍스트를 토큰 리스트로 분할 (자식 클래스에서 구현)"""

    def encode(self, text: str, add_bos=False, add_eos=False) -> list[int]:
        tokens = self._tokenize(text)
        ids = [self.token_to_id.get(t, self.special.unk_id) for t in tokens]
        if add_bos: ids = [self.special.bos_id] + ids
        if add_eos: ids = ids + [self.special.eos_id]
        return ids

    def decode(self, ids: list[int]) -> str:
        tokens = [self.id_to_token[i] for i in ids
                  if self.id_to_token[i] not in self.special.tokens]
        raw = self._tokens_to_bytes(tokens)
        return self._postprocess(raw.decode("utf-8", errors="replace"))
\end{lstlisting}

\begin{infobox}[추상 클래스의 역할]
\code{Tokenizer}는 ``인터페이스''를 정의한다. 자식 클래스(BPE, Unigram)는 \code{\_train}과 \code{\_tokenize}를 구현해야 한다. 이렇게 하면 학습 알고리즘은 달라도 \code{encode/decode} API는 동일하게 유지되어, 상위 코드(데이터셋, 트레이너)가 토크나이저 종류에 무관하게 동작한다.

이것이 ``다형성(polymorphism)''의 가치다. 트레이너는 \code{tokenizer.encode(text)}만 호출하면 되고, 내부적으로 BPE가 돌든 Unigram이 돌든 신경 쓰지 않는다. 나중에 ``SentencePiece''나 ``WordPiece'' 같은 새 알고리즘을 추가해도 트레이너 코드는 수정할 필요가 없다.
\end{infobox}

\section{BPE 토크나이저}

\keyterm{BPE(Byte Pair Encoding)}는 1994년 데이터 압축을 위해 제안된 알고리즘을 2016년 Sennrich et al.이 신경망 기계번역에 도입한 것이다. GPT-2, GPT-3, LLaMA 등이 사용한다.

\subsection{알고리즘 직관}

BPE의 핵심 아이디어는 단순하다:
\begin{enumerate}
  \item 텍스트를 바이트 단위로 분해한다 (모든 문자를 표현 가능, OOV 없음).
  \item 인접한 바이트 쌍 중 가장 자주 등장하는 쌍을 찾는다.
  \item 그 쌍을 새로운 토큰으로 병합한다.
  \item 목표 어휘 크기에 도달할 때까지 2$\sim$3을 반복한다.
\end{enumerate}

처음에는 모든 텍스트가 ``바이트의 나열''이다. 예를 들어 ``lower''는 \code{[l, o, w, e, r]}로 표현된다. 학습을 진행하면서:
\begin{itemize}
  \item ``lo''가 자주 등장하면 병합: \code{[lo, w, e, r]}
  \item ``er''가 자주 등장하면 병합: \code{[lo, w, er]}
  \item ``low''가 자주 등장하면 병합: \code{[low, er]}
  \item ``lower''가 자주 등장하면 병합: \code{[lower]}
\end{itemize}

최종적으로 ``lower''가 하나의 토큰이 되지만, 희귀한 단어는 부분단어로 남는다. ``lowering''은 ``lower'' + ``ing''이 될 수 있다.

\begin{figure}[htbp]
\centering
\begin{tikzpicture}[
  every node/.style={font=\small\sffamily},
  byte/.style={draw=inkblue, fill=inkblue!8, minimum width=0.8cm, minimum height=0.6cm, rounded corners=2pt},
  merge/.style={draw=inkred, fill=inkred!10, minimum width=1.6cm, minimum height=0.6cm, rounded corners=2pt},
  arr/.style={-{Stealth[length=4pt]}, inkgray}
]
  \node[byte] (a1) at (0, 0) {l};
  \node[byte, right=0.1cm of a1] (a2) {o};
  \node[byte, right=0.1cm of a2] (a3) {w};
  \node[byte, right=0.1cm of a3] (a4) {e};
  \node[byte, right=0.1cm of a4] (a5) {r};
  \node[inkgray, anchor=west, font=\scriptsize] at (5.5, 0) {step 0: bytes};

  \node[merge] (b1) at (0, -1) {lo};
  \node[byte, right=0.1cm of b1] (b2) {w};
  \node[byte, right=0.1cm of b2] (b3) {e};
  \node[byte, right=0.1cm of b3] (b4) {r};
  \node[inkgray, anchor=west, font=\scriptsize] at (5.5, -1) {step 1: merge 'l'+'o'};

  \node[merge] (c1) at (0, -2) {low};
  \node[byte, right=0.1cm of c1] (c2) {e};
  \node[byte, right=0.1cm of c2] (c3) {r};
  \node[inkgray, anchor=west, font=\scriptsize] at (5.5, -2) {step 2: merge 'lo'+'w'};

  \node[merge] (d1) at (0, -3) {low};
  \node[merge, right=0.1cm of d1] (d2) {er};
  \node[inkgray, anchor=west, font=\scriptsize] at (5.5, -3) {step 3: merge 'e'+'r'};

  \node[merge, fill=inkblue!20, draw=inkblue, line width=0.8pt] (e1) at (0, -4) {lower};
  \node[inkgray, anchor=west, font=\scriptsize] at (5.5, -4) {step 4: merge 'low'+'er'};

  \foreach \y in {0, -1, -2, -3} {
    \draw[arr] (\y+0.05, \y+0.3) -- (\y+0.05, \y+0.7);
  }
\end{tikzpicture}
\caption{BPE 학습 과정. ``lower''가 바이트에서 시작해 점진적으로 병합되어 하나의 토큰이 되는 과정.}
\label{fig:bpe-process}
\end{figure}

\subsection{왜 바이트 단위인가?}

GPT-2부터 BPE는 ``바이트 단위''로 동작한다. 이전에는 문자(character) 단위를 썼는데, 문제가 있었다:
\begin{itemize}
  \item 이모지, 희귀 한자 등 어휘에 없는 문자는 \code{<unk>} 처리됨.
  \item 다국어 코퍼스에서 문자 집합이 너무 커짐.
\end{itemize}

바이트 단위면 모든 UTF-8 문자를 1$\sim$4바이트로 표현할 수 있고, 어휘의 ``바이트'' 부분은 항상 256개로 고정된다. OOV가 사실상 사라진다.

\subsection{구현: 학습}

\begin{lstlisting}[language=Python]
from collections import Counter
from .base import Tokenizer

class BPETokenizer(Tokenizer):
    def __init__(self, special_tokens=None):
        super().__init__(special_tokens=special_tokens)
        self.merges: list[tuple[str, str]] = []
        self._merge_rank: dict[tuple[str, str], int] = {}

    def _train(self, corpus: list[str], vocab_size: int) -> None:
        # 1. 특수 토큰 + 256개 바이트로 초기 어휘 구성
        self.id_to_token = list(self.special.tokens)
        for b in range(256):
            self.id_to_token.append(chr(b))
        self.token_to_id = {t: i for i, t in enumerate(self.id_to_token)}

        # 2. 코퍼스를 단어 단위로 분할
        word_freqs = Counter()
        for text in corpus:
            for word in self._pre_tokenize(text):
                word_freqs[word] += 1

        # 3. 각 단어를 바이트 토큰 리스트로 변환
        word_to_tokens = {
            w: [chr(b) for b in w.encode("utf-8")]
            for w in word_freqs
        }

        # 4. 병합 루프
        self.merges = []
        target = vocab_size - len(self.id_to_token)
        for _ in range(max(target, 0)):
            pair_freqs = Counter()
            for word, freq in word_freqs.items():
                toks = word_to_tokens[word]
                for i in range(len(toks) - 1):
                    pair_freqs[(toks[i], toks[i+1])] += freq
            if not pair_freqs:
                break
            best = max(pair_freqs.items(), key=lambda x: (x[1], x[0]))
            if best[1] < 2:
                break
            new_token = best[0][0] + best[0][1]
            self.merges.append(best[0])
            self.id_to_token.append(new_token)
            self.token_to_id[new_token] = len(self.id_to_token) - 1
            for word in word_to_tokens:
                word_to_tokens[word] = self._merge_in_list(
                    word_to_tokens[word], best[0], new_token
                )

        self._merge_rank = {p: i for i, p in enumerate(self.merges)}
\end{lstlisting}

\begin{notebox}[병합 우선순위 안정성]
\code{max(pair\_freqs.items(), key=lambda x: (x[1], x[0]))}에서 \code{(x[1], x[0])} 튜플로 정렬하는 이유는 동점 처리 때문이다. 빈도가 같은 쌍이 여러 개일 때, 사전순으로 가장 앞선 쌍을 선택하여 결과를 결정론적으로 만든다. 이렇게 해야 같은 코퍼스로 학습하면 항상 같은 병합 순서가 나온다.
\end{notebox}

\subsection{인코딩}

학습된 병합 규칙을 새 텍스트에 적용하려면, 가장 우선순위가 높은(먼저 학습된) 병합부터 적용한다.

\begin{lstlisting}[language=Python]
    def _tokenize(self, text: str) -> list[str]:
        tokens = []
        for word in self._pre_tokenize(text):
            bs = word.encode("utf-8")
            word_tokens = [chr(b) for b in bs]
            word_tokens = self._apply_merges(word_tokens)
            tokens.extend(word_tokens)
        return tokens

    def _apply_merges(self, tokens: list[str]) -> list[str]:
        """학습된 병합 규칙을 우선순위 순으로 적용."""
        while len(tokens) >= 2:
            # 가장 rank가 낮은 (우선순위 높은) 쌍 찾기
            best_idx, best_rank = -1, len(self.merges) + 1
            for i in range(len(tokens) - 1):
                rank = self._merge_rank.get((tokens[i], tokens[i+1]), -1)
                if 0 <= rank < best_rank:
                    best_rank, best_idx = rank, i
            if best_idx < 0:
                break
            merged = tokens[best_idx] + tokens[best_idx + 1]
            tokens = tokens[:best_idx] + [merged] + tokens[best_idx + 2:]
        return tokens
\end{lstlisting}

\subsection{전처리: 공백 처리}

BPE는 공백을 특별히 처리해야 단어 경계를 보존할 수 있다. GPT-2는 공백을 특수 문자 Ġ(U+0120)로 치환한다. 우리는 SentencePiece 방식을 따라 U+2581(▁)로 치환한다.

\begin{lstlisting}[language=Python]
    def _pre_tokenize(self, text: str) -> list[str]:
        text = text.replace(" ", "▁")
        text = text.replace("\n", "▁").replace("\t", "▁")
        while "▁▁" in text:
            text = text.replace("▁▁", "▁")
        if not text.startswith("▁"):
            text = "▁" + text
        parts = text.split("▁")
        return ["▁" + p for i, p in enumerate(parts) if i > 0 and p]
\end{lstlisting}

\begin{warnbox}[공백 치환의 함정]
``Hello world''를 ``▁Hello▁world''로 치환한 후 split하면 \code{["▁Hello", "▁world"]}가 된다. 단어마다 ▁이 앞에 붙어 있어, 디코딩 시 ▁을 다시 공백으로 바꾸면 원래 문장이 복원된다.

하지만 ``Hello, world!''처럼 구두점이 있으면 주의해야 한다. 단순 split은 ``▁Hello,''를 만들어 쉼표가 단어의 일부가 되어버린다. 실제 GPT-2는 정규식 패턴으로 단어/구두점을 분리하지만, 이 책에서는 단순화를 위해 생략한다.
\end{warnbox}

\subsection{디코딩: 바이트 복원}

BPE의 까다로운 부분은 디코딩이다. 토큰이 \code{chr(byte)}들의 시퀀스이므로, 각 문자의 코드포인트를 바이트로 취급하여 원래 바이트 스트림을 복원한 뒤 UTF-8로 디코딩해야 한다.

\begin{lstlisting}[language=Python]
    def _tokens_to_bytes(self, tokens: list[str]) -> bytes:
        """BPE 토큰을 바이트로 변환. 각 문자의 코드포인트를 바이트로."""
        out = bytearray()
        for tok in tokens:
            for c in tok:
                out.append(ord(c) & 0xFF)
        return bytes(out)
\end{lstlisting}

예를 들어 ``▁Hello'' 토큰이 있으면, ▁(U+2581)은 UTF-8로 \code{0xE2 0x96 0x81}이므로 3바이트이지만, BPE 학습 시 ``▁'' 문자 자체를 3개의 개별 바이트 토큰으로 처리했으므로 \code{chr(0xE2) + chr(0x96) + chr(0x81)}로 표현된다. 디코딩 시 각 \code{chr(b)}에서 \code{ord()}로 바이트를 추출하면 원래 바이트 스트림 \code{0xE2 0x96 0x81 0x48 0x65 0x6C 0x6C 0x6F}를 얻고, 이를 UTF-8로 디코딩하면 ``▁Hello''가 된다.

\section{Unigram 토크나이저}

\keyterm{Unigram} 토크나이저는 Kudo (2018)가 SentencePiece의 일부로 제안한 방식이다. BPE가 ``빈도 기반 병합''이라면, Unigram은 ``확률 기반 분할''이다.

\subsection{직관}

Unigram은 각 토큰에 확률을 할당한다. 문장 ``lower''에 대해 가능한 분할은 여러 가지다:
\begin{itemize}
  \item \code{[l, o, w, e, r]}: 5개 토큰
  \item \code{[lo, w, er]}: 3개 토큰
  \item \code{[low, er]}: 2개 토큰
  \item \code{[lower]}: 1개 토큰
\end{itemize}

각 분할의 확률은 토큰 확률의 곱으로 계산된다. 가장 높은 확률을 갖는 분할을 선택한다. ``lower''가 어휘에 있고 확률이 높으면 \code{[lower]}가 선택되지만, 없으면 더 작은 단위로 분할된다.

\subsection{알고리즘}

\begin{enumerate}
  \item 코퍼스에서 모든 부분문자열을 후보 토큰으로 추출하고 빈도를 센다.
  \item 빈도를 확률로 변환하여 각 토큰에 로그 확률 점수를 할당.
  \item EM(Expectation-Maximization) 반복:
    \begin{itemize}
      \item E-step: Viterbi 알고리즘으로 각 단어의 최적 분할 탐색.
      \item M-step: 분할 결과로 토큰 확률 갱신.
    \end{itemize}
  \item 가장 점수가 낮은 토큰을 제거하며 어휘 축소.
\end{enumerate}

\subsection{Viterbi 디코딩}

주어진 단어에 대해 ``가장 높은 총 로그 확률을 갖는 토큰 분할''을 동적 계획법으로 찾는다.

\begin{lstlisting}[language=Python]
    def _viterbi(self, word: str) -> tuple[list[str], float]:
        n = len(word)
        # dp[i] = (최대 점수, 이전 인덱스, 이전 토큰)
        dp = [(-float("inf"), -1, "") for _ in range(n + 1)]
        dp[0] = (0.0, -1, "")
        for i in range(1, n + 1):
            for j in range(max(0, i - 16), i):
                substr = word[j:i]
                score = self.token_scores.get(substr, -20.0)
                cand = dp[j][0] + score
                if cand > dp[i][0]:
                    dp[i] = (cand, j, substr)
        # 역추적
        path, i = [], n
        while i > 0:
            _, j, tok = dp[i]
            path.append(tok)
            i = j
        path.reverse()
        return path, dp[n][0]
\end{lstlisting}

\begin{infobox}[Viterbi의 복잡도]
$O(n^2)$ 대신 $O(n \cdot L)$로 줄인 비결은 \code{max(0, i-16)} 부분이다. 토큰 길이가 16을 넘는 경우가 드물기 때문에, $j$의 범위를 제한하여 탐색 공간을 줄인다. 실제 SentencePiece도 비슷한 최적화를 사용한다.
\end{infobox}

\section{BPE vs Unigram 비교}

\begin{center}
\begin{tabularx}{\textwidth}{lXX}
\toprule
 & \textbf{BPE} & \textbf{Unigram} \\
\midrule
학습 방식 & 빈도 기반 병합 & 확률 기반 분할 (EM) \\
인코딩 & 병합 규칙 순서 적용 & Viterbi로 최적 분할 \\
결정론적 & O (같은 입력, 같은 출력) & O \\
다중 분할 지원 & X & O (샘플링 가능) \\
주요 사용처 & GPT-2, GPT-3, LLaMA & T5, ALBERT, mBART \\
\bottomrule
\end{tabularx}
\end{center}

\section{토크나이저 테스트}

구현이 끝났으면 단위 테스트로 검증한다. 가장 중요한 테스트는 \keyterm{라운드트립(round-trip) 테스트}: 인코딩 후 디코딩하면 원문이 복원되어야 한다는 것이다.

\begin{lstlisting}[language=Python]
def test_encode_decode_roundtrip():
    tok = BPETokenizer()
    tok.train(CORPUS, vocab_size=200)
    text = "the quick fox"
    ids = tok.encode(text)
    decoded = tok.decode(ids)
    assert " ".join(decoded.split()) == "the quick fox"

def test_bos_eos_addition():
    tok = BPETokenizer()
    tok.train(CORPUS, vocab_size=100)
    ids = tok.encode("hello", add_bos=True, add_eos=True)
    assert ids[0] == tok.special.bos_id
    assert ids[-1] == tok.special.eos_id

def test_unknown_token_handled():
    """어휘에 없는 문자도 바이트 단위로 처리되어 unk가 발생하지 않아야 함."""
    tok = BPETokenizer()
    tok.train(CORPUS, vocab_size=100)
    ids = tok.encode("안녕")
    assert tok.special.unk_id not in ids  # 바이트로 분해되므로 unk 없음

def test_save_load(tmp_path):
    tok = BPETokenizer()
    tok.train(CORPUS, vocab_size=150)
    path = tmp_path / "tok.json"
    tok.save(path)
    tok2 = BPETokenizer.load(path)
    text = "the quick fox"
    assert tok.encode(text) == tok2.encode(text)
\end{lstlisting}

\begin{notebox}[테스트 철학]
이 책의 모든 모듈은 세 가지 수준의 테스트를 갖는다:
\begin{itemize}
  \item \textbf{형상 테스트(shape test)}: 텐서의 shape이 예상과 일치하는지.
  \item \textbf{수치 테스트(numeric test)}: 작은 입력에 대해 손계산 결과와 비교.
  \item \textbf{통합 테스트}: 여러 모듈을 연결했을 때 end-to-end로 동작.
\end{itemize}
이 세 가지를 통과한 코드만 책에 수록된다.

라운드트립 테스트는 통합 테스트의 한 예다. ``인코딩 후 디코딩하면 원문이 나온다''는 자명해 보이지만, 실제로는 공백 처리, 멀티바이트 문자, 특수 토큰 제거 등 여러 단계에서 버그가 생길 수 있다. 이 테스트가 통과하면 전체 파이프라인이 올바르게 동작함을 확신할 수 있다.
\end{notebox}

\section{실습: 실제 말뭉치로 토크나이저 학습}

책의 \code{scripts/tiny\_train.py} 데모를 실행하면 실제로 BPE 토크나이저를 학습하는 과정을 볼 수 있다.

\begin{lstlisting}[style=bashstyle]
cd make-llm-basic
python scripts/tiny_train.py --vocab 500
\end{lstlisting}

출력 예:
\begin{lstlisting}[basicstyle=\ttfamily\footnotesize, backgroundcolor=\color{codebg}, frame=single]
[1/5] 토크나이저 학습 중...
    완료: 308개 토큰, 1.2초
    샘플: 'the quick brown fox jumps over the lazy dog'
    → 18개 토큰: [264, 282, 291, 267, 261, 110, ...]
\end{lstlisting}

\section{요약}

\begin{itemize}
  \item 토크나이저는 텍스트를 정수 시퀀스로 변환하는 LLM의 첫 단계다.
  \item 서브워드 방식이 어휘 크기와 OOV 사이의 균형을 잡는다.
  \item BPE는 빈도 기반 병합, Unigram은 확률 기반 분할.
  \item 특수 토큰 \code{<pad>/<bos>/<eos>/<unk>}는 ID 0$\sim$3에 고정.
  \item 바이트 단위 BPE는 OOV를 사실상 제거한다.
  \item Viterbi 알고리즘으로 Unigram의 최적 분할을 $O(n)$에 찾는다.
  \item 모든 구현은 pytest 단위 테스트로 검증된다.
\end{itemize}

\section{연습 문제}

\begin{enumerate}
  \item BPE 학습 시 ``가장 빈도가 높은 쌍'' 대신 ``가장 빈도가 낮은 쌍''을 병합하면 어떻게 될까? 결과를 예측하고 실험해보라.
  \item Unigram 토크나이저에서 \code{score = -20.0} 패널티 값을 \code{-5.0}으로 바꾸면 어떤 일이 일어나는가?
  \item 한국어 코퍼스 ``안녕하세요 만나서 반갑습니다''를 BPE로 학습한다고 하자. 어휘 크기가 100일 때와 1000일 때 어떤 차이가 나는가?
  \item Viterbi에서 토큰 길이 상한을 16이 아니라 4로 줄이면 어떤 토큰이 사라지는가?
  \item \code{<bos>} 토큰이 왜 필요한지 생각해보라. 없어도 모델이 동작할까?
\end{enumerate}

다음 장에서는 토큰 ID를 밀집 벡터로 변환하는 임베딩 레이어를 구현한다.
